\chapter{Reproducibility and Implementation Details}
\label{app:reproducibility}

This appendix records the implementation details needed to understand and reproduce the experiments in this dissertation. The appendix therefore summarises the model configuration, prompt templates, evaluation pipeline, output formats, metric definitions, generated files, and validation checks without reproducing the full codebase.

\section{Experimental Configuration}
\label{app:experimental-configuration}

Table~\ref{tab:appendix_config} summarises the fixed configuration used for the LLM calls and evaluation. The model settings were read from \texttt{configs/model.yaml}; these values match the configuration requested for the dissertation experiments.

\begin{table}[htbp]
    \centering
    \caption{Experimental configuration used for the main prompt-comparison pipeline}
    \label{tab:appendix_config}
    \renewcommand{\arraystretch}{1.15}
    \begin{adjustbox}{max width=\textwidth}
    \begin{tabular}{ll}
        \toprule
        \textbf{Item} & \textbf{Value} \\
        \midrule
        Model & OpenAI GPT-4o \\
        Temperature & 0.0 \\
        Top-p & 1.0 \\
        Maximum output tokens & 128 \\
        Seed & 1234 \\
        Request timeout & 60 seconds \\
        Dataset & W\&I + LOCNESS, BEA-2019 shared task \\
        Main experiment size & 500 learner sentences \\
        Sentence-length experiment size & 450 learner sentences \\
        Evaluation tool & ERRANT \\
        Main accuracy metric & ERRANT $F_{0.5}$ \\
        Main over-correction indicator & Composite Over-Correction Index (OCI) \\
        \bottomrule
    \end{tabular}
    \end{adjustbox}
\end{table}

The main comparison used the same 500 learner sentences for each prompt condition. The sentence-length analysis used 450 learner sentences, with 150 short, 150 medium, and 150 long sentences. The scripts use the placeholder \texttt{<SENTENCE>} in prompt templates and replace it with the learner sentence before sending the request to the model.

The dataset preparation first produced a reproducible pool of candidate sentences from the processed W\&I + LOCNESS data using random seed 42. From this pool, the main prompt-refinement and prompt-comparison experiments used a fixed 500-sentence evaluation set. This ensured that all prompt conditions were evaluated on identical inputs.

For the sentence-length experiment, a balanced subset was required so that short, medium, and long sentences could be compared fairly. Sentences were assigned to length groups using source-token counts, and 150 sentence IDs were selected for each group using seed 1234. Where possible, these sentences were drawn from the main 500-sentence evaluation set. The final sentence-length subset contained 450 unique sentence IDs: 417 overlapped with the main 500-sentence set, while 33 additional examples were selected from the same candidate pool to complete the balanced length groups.

This means that the sentence-length analysis was not based on a separate dataset, but on a controlled balanced subset derived from the same prepared corpus pipeline.

\section{Full Prompt Templates}
\label{app:prompt-templates}

The following prompt templates are reproduced exactly from the project prompt files. The final selected prompts after Experiment 1 were Instruction v2, Role v4, and Few-shot v4. In Experiments 2--6 these are referred to as Instruction, Role, and Few-shot respectively.

\subsection{Baseline prompt}
\begin{verbatim}
Please correct any grammatical errors in this sentence:
<SENTENCE>
Return only the corrected sentence.
\end{verbatim}

\subsection{Instruction v1}
\begin{verbatim}
Correct the grammatical errors in the following sentence:
<SENTENCE>
Return only the corrected sentence. No explanation.
\end{verbatim}

\subsection{Instruction v2 (final selected instruction prompt)}
\textbf{Selected prompt for the main comparison.}

\begin{verbatim}
Correct only grammatical errors in the sentence below.
Do not change word choice or sentence structure.
Return only the corrected sentence. No explanation.
<SENTENCE>
\end{verbatim}

\subsection{Instruction v3}
\begin{verbatim}
Correct only grammatical errors in the sentence below.
Do not change word choice, tone, or sentence structure unless strictly required for grammar.
Keep the original meaning and style intact.
Return only the corrected sentence. No explanation.
<SENTENCE>
\end{verbatim}

\subsection{Instruction v4}
\begin{verbatim}
Correct only grammatical errors in the sentence below.
Do not change word choice, tone, or sentence structure unless strictly required for grammar.
Do not add or remove words unnecessarily.
Do not improve style or vocabulary.
Keep the original meaning and style intact.
Return only the corrected sentence. No explanation.
<SENTENCE>
\end{verbatim}

\subsection{Role v1}
\begin{verbatim}
You are an English teacher.
Correct the grammatical errors in the following sentence:
<SENTENCE>
Return only the corrected sentence. No explanation.
\end{verbatim}

\subsection{Role v2}
\begin{verbatim}
You are an English teacher correcting an L2 student's sentence.
Make only grammatical corrections.
Return only the corrected sentence. No explanation.
<SENTENCE>
\end{verbatim}

\subsection{Role v3}
\begin{verbatim}
You are an English teacher correcting an L2 student's sentence.
Make minimal grammatical corrections only. Do not improve style.
Do not change word choice or sentence structure.
Return only the corrected sentence. No explanation.
<SENTENCE>
\end{verbatim}

\subsection{Role v4 (final selected role-based prompt)}
\textbf{Selected prompt for the main comparison.}

\begin{verbatim}
You are an English teacher correcting an L2 student's sentence.
Make only the minimal grammatical corrections needed.
Do not change word choice, tone, or sentence structure.
Do not add or remove words unnecessarily.
Do not improve vocabulary or style.
Return only the corrected sentence. No explanation.
<SENTENCE>
\end{verbatim}

\subsection{Few-shot v1}
\begin{verbatim}
Correct the grammatical errors in the following sentence.

Example:
Input: She go to university every day.
Output: She goes to university every day.

Now correct this sentence:
<SENTENCE>
Return only the corrected sentence. No explanation.
\end{verbatim}

\subsection{Few-shot v2}
\begin{verbatim}
Correct the grammatical errors in the following sentence.

Example 1:
Input: She go to university every day.
Output: She goes to university every day.

Example 2:
Input: I look forward to meet you.
Output: I look forward to meeting you.

Now correct this sentence:
<SENTENCE>
Return only the corrected sentence. No explanation.
\end{verbatim}

\subsection{Few-shot v3}
\begin{verbatim}
Correct the grammatical errors in the following sentence.

Example 1:
Input: She go to university every day.
Output: She goes to university every day.

Example 2:
Input: I look forward to meet you.
Output: I look forward to meeting you.

Example 3:
Input: He bought car yesterday.
Output: He bought a car yesterday.

Now correct this sentence:
<SENTENCE>
Return only the corrected sentence. No explanation.
\end{verbatim}

\subsection{Few-shot v4 (final selected few-shot prompt)}
\textbf{Selected prompt for the main comparison.}

\begin{verbatim}
You are correcting sentences written by an L2 learner.
Correct only grammatical errors. Keep the original wording and structure whenever possible.

Example 1:
Input: She go to university every day.
Output: She goes to university every day.

Example 2:
Input: I look forward to meet you.
Output: I look forward to meeting you.

Example 3:
Input: He bought car yesterday.
Output: He bought a car yesterday.

Example 4:
Input: The book is on table.
Output: The book is on the table.

Now correct this sentence:
<SENTENCE>
Return only the corrected sentence. No explanation.
\end{verbatim}

\section{Example Output Record Format}
\label{app:output-record-format}

Model outputs were stored as JSONL records, with one JSON object per sentence--prompt pair. The real records also include model metadata such as temperature, top-p, seed, timeout, attempt number, elapsed time, raw model output, and cleaned model output. The example below is a representative non-sensitive record from the project, shortened to the fields most relevant for reproducibility.

\begin{verbatim}
{
  "sentence_id": "B.train.gold.bea19_001618",
  "prompt_type": "fewshot",
  "prompt_version": "selected final prompt",
  "source": "So , it would be a good choice to consider them your friends .",
  "reference": "So , it would be a good choice to consider them your friends .",
  "model_output": "So, it would be a good choice to consider them your friends."
}
\end{verbatim}

\section{Metric Implementation and OCI Details}
\label{app:metric-implementation}

The evaluation used ERRANT precision, recall, and $F_{0.5}$ for reference-based grammatical correction accuracy. OCI was then used as a comparative indicator of over-correction risk. It should not be read as a definitive judgement that an individual edit is unnecessary. Instead, it estimates how far the model output moves away from the source after accounting for estimated utility.

The component metrics were implemented as follows:
\begin{itemize}
    \item \textbf{Word-level edit distance}: Levenshtein distance between whitespace-tokenised source and output sentences.
    \item \textbf{Edit density}: word-level edit distance divided by the source word count, using a denominator of at least 1.
    \item \textbf{Delta TTR}: the absolute difference between source and output type-token ratio after lower-casing and punctuation stripping.
    \item \textbf{Delta readability}: the absolute difference between source and output Flesch--Kincaid grade estimates.
    \item \textbf{Cosine similarity}: cosine similarity between lightweight stylometric feature vectors for the source and output. The feature vector includes sentence length, average word length, punctuation count, and simple lexical proportion features.
    \item \textbf{Delta fluency}: output fluency score minus source fluency score. Fluency is a deterministic surface proxy that penalises adjacent repetition, unusual non-alphabetic characters, high punctuation density, unusual average word length, long sentences, and missing terminal punctuation.
\end{itemize}

For a component value $x$, min--max normalisation was computed as:
\begin{equation}
    \operatorname{norm}(x)=
    \begin{cases}
    0, & x_{\max}=x_{\min} \\
    \dfrac{x-x_{\min}}{x_{\max}-x_{\min}}, & \text{otherwise.}
    \end{cases}
\end{equation}

The divergence component was calculated as:
\begin{align}
OCI_{\text{divergence}} ={}& 0.35 \cdot \operatorname{norm}(\text{edit distance})
+ 0.15 \cdot \operatorname{norm}(\text{edit density}) \\
&+ 0.20 \cdot \operatorname{norm}(\Delta TTR)
+ 0.15 \cdot \operatorname{norm}(\Delta readability) \\
&+ 0.15 \cdot (1 - \text{cosine similarity}).
\end{align}

The utility term was calculated as:
\begin{equation}
\text{Utility} = 0.5 \cdot \operatorname{norm}(\Delta \text{fluency}) + 0.5 \cdot \text{cosine similarity}.
\end{equation}

The final OCI used in the dissertation was:
\begin{equation}
OCI = OCI_{\text{divergence}} \cdot (1 - \text{Utility}).
\end{equation}

Values were clipped to the range $[0,1]$ after the utility and final OCI calculations. Cosine distance, $1 - \text{cosine similarity}$, was also clipped to $[0,1]$. If a min--max denominator was zero, the normalised value was set to 0. Unchanged sentences are therefore handled naturally by the same pipeline: they normally have zero or near-zero edit distance and edit density, and their final OCI depends on the remaining component values after normalisation.

The OCI weights were set using a criterion-based design rationale rather than being learned from human judgement data. Direct edit distance was given the largest weight because it is the most observable signal of source-output rewriting. Edit density was included to account for sentence length, while lexical diversity, readability, and cosine similarity were included to capture broader changes in vocabulary, structure, and source-output similarity. Since these weights involve design judgement, Experiment 4 tests whether the main prompt rankings remain stable under alternative weighting schemes. This means that OCI should be interpreted as a comparative prompt-level indicator rather than a definitive judgement of individual edits.

\section{Generated Files and Results Outputs}
\label{app:generated-files}

Table~\ref{tab:appendix_outputs} summarises the main files produced by the pipeline. The paths are shortened and grouped by output type so that the table shows the structure of the results without listing every individual file.

\begin{table}[htbp]
    \centering
    \caption{Summary of the main output files generated by the experimental pipeline}
    \label{tab:appendix_outputs}
    \renewcommand{\arraystretch}{1.12}
    \begin{adjustbox}{max width=\textwidth}
    \begin{tabular}{p{0.25\textwidth}p{0.34\textwidth}p{0.34\textwidth}}
        \toprule
        \textbf{Output category} & \textbf{Example file/location} & \textbf{Purpose} \\
        \midrule
        Model output records & \texttt{outputs/experiment\_*/*outputs.jsonl} & JSONL records containing model corrections, source/reference text, and run metadata. \\
        ERRANT input files & \texttt{errant\_inputs/*.src, *.hyp, *.ref} & Source, hypothesis, and reference files prepared for ERRANT scoring. \\
        ERRANT reports & \texttt{errant\_outputs/*.report.txt} & Precision, recall, and $F_{0.5}$ reports for each prompt condition. \\
        Style/component metrics & \texttt{style\_eval/*.csv} & Sentence-level and aggregate edit distance, edit density, TTR, readability, fluency, and cosine-similarity components. \\
        OCI results & \texttt{oci\_eval/*.csv} & Sentence-level and aggregate OCI divergence, utility, and final OCI values. \\
        Sentence-length results & \texttt{experiment\_3\_sentence\_length/*.csv} & Length-group outputs and aggregated results for short, medium, and long sentences. \\
        OCI robustness results & \texttt{experiment\_4\_oci\_robustness/*.csv} & OCI rankings and sensitivity results under alternative weighting schemes. \\
        Trade-off results & \texttt{experiment\_5\_tradeoff/*.csv} & Combined ERRANT and OCI tables used for the accuracy--over-correction trade-off analysis. \\
        Qualitative examples & \texttt{experiment\_6\_qualitative/*.csv} & Selected examples used for sentence-level qualitative discussion. \\
        \bottomrule
    \end{tabular}
    \end{adjustbox}
\end{table}

\section{Pipeline Workflow}
\label{app:pipeline-workflow}

The pipeline followed the same broad workflow across the experiments:
\begin{enumerate}
    \item Load the prepared W\&I + LOCNESS sentence CSV containing sentence IDs, source sentences, and reference corrections.
    \item Load the relevant prompt templates and render each template by replacing \texttt{<SENTENCE>} with the learner sentence.
    \item Call GPT-4o using the fixed configuration in Table~\ref{tab:appendix_config}.
    \item Store raw and cleaned model outputs as JSONL records with metadata.
    \item Convert JSONL outputs into ERRANT \texttt{.src}, \texttt{.hyp}, and \texttt{.ref} files.
    \item Run ERRANT to obtain precision, recall, and $F_{0.5}$.
    \item Compute sentence-level style and surface-change components.
    \item Compute sentence-level and aggregate OCI values.
    \item Combine ERRANT and OCI results for prompt comparison, robustness analysis, trade-off analysis, and qualitative inspection.
\end{enumerate}

The model-output cleaning step was deterministic. It trimmed whitespace, kept the first non-empty output line, removed wrapping quotes or backticks, removed leading labels such as ``Corrected:'' or ``Output:'', and collapsed repeated whitespace. This was used to reduce formatting noise before ERRANT and OCI evaluation.

\section{Testing and Validation Checks}
\label{app:testing-validation}

The project included several validation checks in the data preparation and evaluation scripts. Where explicit log files were not retained, the checks below describe the validation logic implemented in the scripts and the current output evidence available in the project folder.

\begin{table}[htbp]
    \centering
    \caption{Line-count validation for Experiment 2 ERRANT input files}
    \label{tab:appendix_line_counts}
    \renewcommand{\arraystretch}{1.12}
    \begin{adjustbox}{max width=\textwidth}
    \begin{tabular}{lrrrr}
        \toprule
        \textbf{Condition} & \textbf{.src lines} & \textbf{.hyp lines} & \textbf{.ref lines} & \textbf{Status} \\
        \midrule
        Experiment 2 baseline & 500 & 500 & 500 & match \\
        Experiment 2 instruction & 500 & 500 & 500 & match \\
        Experiment 2 role & 500 & 500 & 500 & match \\
        Experiment 2 few-shot & 500 & 500 & 500 & match \\
        \bottomrule
    \end{tabular}
    \end{adjustbox}
\end{table}

Additional checks used in the implementation were:
\begin{itemize}
    \item The current Experiment 2 JSONL file contains 2000 non-empty records. Sentence ID coverage by prompt is: baseline: 500, fewshot: 500, instruction: 500, role: 500.
    \item Empty cleaned outputs in the current Experiment 2 JSONL file: 0. Malformed JSONL lines detected during this check: 0.
    \item ERRANT preparation scripts skip malformed JSON rows and rows with empty source, hypothesis, or reference fields, and they raise an error if \texttt{.src}, \texttt{.hyp}, and \texttt{.ref} line counts do not match.
    \item Sentence-length selection uses a fixed random seed and checks the number of selected sentence IDs in each length group before generating or consolidating outputs.
    \item Resume logic records completed \texttt{(sentence\_id, prompt\_id)} pairs, which reduces the risk of duplicate LLM calls during interrupted runs.
    \item Metric sanity checks include safe division for zero denominators, min--max handling when all values are identical, clipping OCI-related values to $[0,1]$, and aggregate summaries at both sentence and prompt level.
    \item Prompt selection was validated using both ERRANT $F_{0.5}$ and OCI: Experiment 1 selected the strongest representative from each prompt family before the main strategy comparison.
\end{itemize}